\documentclass[11pt,a4paper]{article}

% Appunti personali di lettura: non sono un testo da presentare.
\usepackage[utf8]{inputenc}
\usepackage[T1]{fontenc}
\usepackage[italian]{babel}
\usepackage{amsmath,amssymb,mathtools}
\usepackage[a4paper,margin=2.45cm]{geometry}
\usepackage{enumitem}
\usepackage{booktabs}
\usepackage{xcolor}
\usepackage{tikz}
\usetikzlibrary{arrows.meta,positioning,calc,patterns}
\usepackage{hyperref}

\title{Appunti di lettura\\\large Schmidt (2005), ``Tail Dependence''}
\author{Appunti personali di Massimo}
\date{}

\newcommand{\R}{\mathbb{R}}
\newcommand{\Prob}{\mathbb{P}}
\newcommand{\ind}{\mathbf{1}}
\newcommand{\cdf}{\operatorname{F}}
\newcommand{\Idea}{\paragraph{Idea.}}
\newcommand{\Ricorda}{\paragraph{Da ricordare.}}
\newcommand{\Esempio}{\paragraph{Esempio.}}

\begin{document}
\maketitle

\paragraph{Come usare questi appunti.}
Questo documento serve a ricostruire l'idea del capitolo prima di discuterla
con il relatore. Le frasi introdotte da ``Nel paper'' riassumono ciò che
Schmidt afferma; i paragrafi ``Idea'', ``Esempio'' e ``Da ricordare'' traducono
la matematica in parole. Le figure aggiunte qui sono diagrammi concettuali,
non ricostruzioni dei grafici o dei dati originali del paper. I numeri tra
parentesi rimandano alle sezioni e alle formule del paper (pp.~65--91 del
volume).

\section{Perché parlare di tail dependence? (Sec.~3.1, pp.~65--66)}

La tail dependence misura quanto due variabili continuino a muoversi insieme
quando ci si sposta verso una coda della distribuzione: in basso a sinistra o
in alto a destra. Se $X_1$ e $X_2$ sono rendimenti, il basso a sinistra
rappresenta rendimenti molto negativi, cioè grandi perdite; se invece si
trasformano i rendimenti in perdite positive, le stesse osservazioni finiscono
nella coda superiore, in alto a destra. L'interesse è soprattutto finanziario:
il Value at Risk (VaR)
di un portafoglio dipende sia dalle code marginali dei singoli titoli sia da
come gli eventi estremi si presentano insieme.

Nel paper la correlazione lineare è presentata come misura spesso
insufficiente. Può descrivere bene la dipendenza media, ma non dice quanto sia
probabile una perdita estrema di un titolo quando un'altra perdita estrema è
già avvenuta. Una specifica errata della dipendenza di coda può quindi rendere
inaffidabile il VaR.

L'obiettivo dichiarato è introdurre il tail-dependence coefficient (TDC),
calcolarlo per varie famiglie di distribuzioni e copule, discutere tre
stimatori (due non parametrici e uno parametrico), confrontarli in simulazione
e mostrare un'applicazione al VaR.

\begin{figure}[ht]
\centering
\begin{tikzpicture}[x=0.40cm,y=0.40cm,font=\small]
  % Primo pannello: punti con la stessa relazione centrale, ma senza cluster estremo.
  \begin{scope}
    \draw[->] (-4.2,0)--(4.4,0) node[right] {$X_1$};
    \draw[->] (0,-4.2)--(0,4.4) node[above] {$X_2$};
    \node[align=center] at (0,6.0) {Normale, nessun cluster estremo};
    \draw[dashed,gray] (-1.45,-4.0)--(-1.45,4.0);
    \draw[dashed,gray] (-4.0,-1.45)--(4.0,-1.45);
    \node[gray,anchor=north east] at (-1.5,-1.5) {coda};
    \foreach \x/\y in {-3.4/-0.2,-2.7/-1.0,-2.1/-1.4,-1.7/-2.7,-1.0/-0.8,-0.4/-1.8,0.3/0.2,0.9/1.0,1.4/1.7,2.2/1.1,2.8/2.6,3.3/2.0,-3.0/1.8,-2.0/2.5,1.8/-1.0,2.9/-0.8}{\fill[blue!55] (\x,\y) circle (1.6pt);}
    \node[align=center] at (0,-5.05) {$\lambda_L\simeq0$};
  \end{scope}
  % Secondo pannello: stesso centro, ma molti punti nella lower tail.
  \begin{scope}[xshift=8cm]
    \draw[->] (-4.2,0)--(4.4,0) node[right] {$X_1$};
    \draw[->] (0,-4.2)--(0,4.4) node[above] {$X_2$};
    \node[align=center] at (0,6.0) {Dipendenza di coda};
    \draw[dashed,gray] (-1.45,-4.0)--(-1.45,4.0);
    \draw[dashed,gray] (-4.0,-1.45)--(4.0,-1.45);
    \node[gray,anchor=north east] at (-1.5,-1.5) {coda};
    \foreach \x/\y in {-3.4/-0.2,-2.7/-1.0,-2.1/-1.4,-1.7/-2.7,-1.0/-0.8,-0.4/-1.8,0.3/0.2,0.9/1.0,1.4/1.7,2.2/1.1,2.8/2.6,3.3/2.0,-3.0/1.8,-2.0/2.5,1.8/-1.0,2.9/-0.8,-3.3/-3.1,-2.9/-2.6,-2.5/-3.2,-2.0/-2.8,-1.5/-3.5,-1.1/-2.9}{\fill[red!65] (\x,\y) circle (1.7pt);}
    \draw[red!70,dashed,thick] (-3.9,-3.9)--(-0.8,-3.9)--(-0.8,-0.8);
    \node[align=center] at (0,-5.05) {$\lambda_L>0$};
  \end{scope}
\end{tikzpicture}
\caption{Diagramma concettuale: due campioni possono avere una relazione centrale simile, ma differire nella concentrazione di osservazioni nella coda inferiore. Non sono dati reali.}
\label{fig:tail-intuition}
\end{figure}

\section{Definizione: una probabilità condizionata in coda (Sec.~3.2, pp.~66--68)}

Sia $X=(X_1,X_2)^\top$ un vettore aleatorio bivariato, con funzioni di
distribuzione marginali $F_1,F_2$ e inverse generalizzate $F_1^{-1},F_2^{-1}$.
Il coefficiente di dipendenza nella coda superiore è
\begin{equation}
  \lambda_U
  =\lim_{v\uparrow 1}
  \Prob\!\left(X_1>F_1^{-1}(v)\mid X_2>F_2^{-1}(v)\right),
  \label{eq:upper}
\end{equation}
quando il limite esiste. Analogamente, nella coda inferiore:
\begin{equation}
  \lambda_L
  =\lim_{v\downarrow 0}
  \Prob\!\left(X_1\leq F_1^{-1}(v)\mid X_2\leq F_2^{-1}(v)\right).
  \label{eq:lower}
\end{equation}

\Idea Fissata una soglia molto alta (o molto bassa), si guarda la frazione di
volte in cui anche la seconda variabile supera la propria soglia, dato che la
prima lo ha già fatto. Il TDC è ciò che rimane quando la soglia viene spinta
all'estremo. È una misura direzionale (upper/lower), non una misura globale
della dipendenza.

Quindi $\lambda_U>0$ (rispettivamente $\lambda_L>0$) significa dipendenza
asintotica nella coda superiore (inferiore); se il coefficiente è zero si
parla di indipendenza asintotica in quella coda. Non si tratta di
indipendenza per ogni soglia: anche quando $\lambda_U=0$, le due variabili
possono essere dipendenti a livelli finiti.

\begin{figure}[ht]
\centering
\begin{tikzpicture}[x=0.78cm,y=0.78cm,font=\small]
  \draw[->] (-4.2,0)--(4.3,0) node[right] {$X_1$};
  \draw[->] (0,-4.2)--(0,4.3) node[above] {$X_2$};
  \draw[dashed,gray] (-1.5,-4)--(-1.5,4) node[above] {$q_1$};
  \draw[dashed,gray] (-4,-1.5)--(4,-1.5) node[right] {$q_2$};
  \fill[orange!20] (-4,-4) rectangle (-1.5,-1.5);
  \node[align=center] at (-2.75,-2.0) {lower-left\\$X_1\le q_1,\ X_2\le q_2$};
  \fill[blue!15] (1.5,1.5) rectangle (4,4);
  \node[align=center] at (2.75,2.2) {upper-right\\$X_1>q_1,\ X_2>q_2$};
  \foreach \x/\y in {-3.5/-3.0,-2.9/-2.4,-2.4/-3.3,-1.9/-2.1,-0.7/-0.5,0.2/0.5,1.1/1.5,1.8/2.0,2.7/2.6,3.3/3.0,-2.8/0.8,0.7/-1.3}{\fill[black] (\x,\y) circle (1.8pt);}
  \draw[->,thick,red!70] (-3.6,-3.55)--(-2.7,-2.75);
  \node[red!70,anchor=north] at (-2.9,-3.55) {una osservazione};
\end{tikzpicture}
\caption{Lettura geometrica: ogni punto è una coppia osservata $(X_1,X_2)$. Le soglie $q_1,q_2$ separano le zone estreme; $\lambda_L$ guarda quanto spesso i punti cadono insieme nella zona lower-left, $\lambda_U$ nella zona upper-right.}
\label{fig:quadrants}
\end{figure}

Per distribuzioni bivariate normali o $t$, il paper usa anche la funzione a
soglia finita
\begin{equation}
  \lambda_U(v)=2\,\Prob\!\left(X_1>F_1^{-1}(v)\mid
  X_2=F_2^{-1}(v)\right),\qquad
  \lambda_U=\lim_{v\uparrow1}\lambda_U(v).
  \label{eq:finite}
\end{equation}
Qui è importante distinguere la curva $\lambda_U(v)$ dal suo limite: la
figura 3.1 o 3.2 non sta dicendo che il TDC sia uguale al valore della curva
per una soglia qualsiasi; mostra come la curva si comporta avvicinandosi a
$v=1$. Se la curva tende a zero, il modello è upper-tail-independent; se tende
a un valore positivo, quel valore è $\lambda_U$. Per le distribuzioni
simmetriche considerate, lower e upper TDC coincidono, ma le figure discutono
esplicitamente la funzione upper $\lambda_U(v)$. La normale bivariata ha
$\lambda_U=0$ per ogni $\rho\in(-1,1)$, mentre la $t$ bivariata ha un limite
positivo, più alto con dipendenza centrale maggiore e con code più pesanti.

\section{Il ruolo delle copule (Sec.~3.2, pp.~67--68)}

\subsection{Tre oggetti da non confondere}
Una osservazione bivariata è una singola coppia di numeri, per esempio
$(r_{\mathrm{BMW}},r_{\mathrm{DB}})$: due rendimenti osservati nello stesso
giorno. La marginale di BMW è il comportamento di $r_{\mathrm{BMW}}$ preso da
solo (e analogamente per Deutsche Bank). Una copula non descrive direttamente
la scala dei rendimenti: descrive come i percentili delle due marginali si
combinano.

\Esempio Se un rendimento è al percentile $0.98$ e l'altro al percentile
$0.97$, la coppia si trova vicino all'angolo upper-right del quadrato dei
percentili, anche se le unità originali dei due rendimenti sono diverse. È
proprio questa posizione relativa nelle code che interessa al TDC.

Una copula $C:[0,1]^2\to[0,1]$ collega le marginali alla distribuzione
congiunta:
\begin{equation}
  F(x_1,x_2)=C\bigl(F_1(x_1),F_2(x_2)\bigr).
  \label{eq:sklar}
\end{equation}

\begin{figure}[ht]
\centering
\begin{tikzpicture}[font=\scriptsize,>=Latex]
  \begin{scope}[xshift=0cm]
    \draw[->] (0,0)--(3.2,0) node[right] {BMW};
    \draw[->] (0,0)--(0,3.2) node[above] {DB};
    \foreach \x/\y in {0.5/0.7,0.8/0.5,1.0/1.2,1.3/1.4,1.5/1.1,1.8/2.0,2.1/1.9,2.4/2.5,2.8/2.3}{\fill[blue!65] (\x,\y) circle (1.5pt);}
    \node[align=center] at (1.5,-0.6) {spazio originale\\marginali in unità diverse};
  \end{scope}
  \draw[->,thick] (3.5,1.5)--(4.5,1.5) node[midway,above] {ranghi};
  \begin{scope}[xshift=5cm]
    \draw[->] (0,0)--(3.2,0) node[right] {$u=F_1(x_1)$};
    \draw[->] (0,0)--(0,3.2) node[above] {$v=F_2(x_2)$};
    \foreach \x/\y in {0.5/0.7,0.8/0.5,1.0/1.2,1.3/1.4,1.5/1.1,1.8/2.0,2.1/1.9,2.4/2.5,2.8/2.3}{\fill[blue!65] ({\x/3},{\y/3}) circle (1.5pt);}
    \node[align=center] at (1.5,-0.6) {percentili\\in $[0,1]^2$};
  \end{scope}
  \draw[->,thick] (8.5,1.5)--(9.5,1.5) node[midway,above] {copula};
  \begin{scope}[xshift=10cm]
    \draw[->] (0,0)--(3.2,0) node[right] {$u$};
    \draw[->] (0,0)--(0,3.2) node[above] {$v$};
    \draw[gray] (0,0) rectangle (3,3);
    \foreach \x/\y in {0.5/0.7,0.8/0.5,1.0/1.2,1.3/1.4,1.5/1.1,1.8/2.0,2.1/1.9,2.4/2.5,2.8/2.3}{\fill[red!65] ({\x/3},{\y/3}) circle (1.5pt);}
    \node[align=center] at (1.5,-0.6) {dipendenza\\tra percentili};
  \end{scope}
\end{tikzpicture}
\caption{Schema concettuale: dai dati originali alle marginali standardizzate e poi alla copula. I pannelli non rappresentano dati reali e le coordinate sono solo illustrative.}
\label{fig:copula-pipeline}
\end{figure}

\paragraph{Modellazione diretta o marginali più copula?}
Modellare direttamente una distribuzione bivariata sceglie in un solo passo
code, scale e dipendenza, ma può imporre simmetrie forti. Separare marginali e
copula permette di usare una distribuzione adatta a ciascun titolo e una
famiglia diversa per il legame tra percentili. La flessibilità è utile per il
VaR, ma introduce due stime e richiede di scegliere la copula giusta.

Per marginali continue il TDC dipende soltanto da $C$, non dalle unità di
misura o dalle trasformazioni strettamente crescenti delle marginali:
\begin{align}
  \lambda_U&=\lim_{v\uparrow1}\frac{1-2v+C(v,v)}{1-v},
  \label{eq:copula-upper}\\
  \lambda_L&=\lim_{v\downarrow0}\frac{C(v,v)}{v}.
  \label{eq:copula-lower}
\end{align}

\section{Famiglie di modelli: quale legame stiamo ipotizzando? (Sec.~3.3)}

Una famiglia non significa dati diversi: significa una forma matematica
ipotizzata per il legame tra le variabili. Il parametro della famiglia cambia
la forza o la forma di quel legame, ma non è automaticamente confrontabile
con un parametro omonimo di un'altra famiglia. Prima di confrontare numeri tra
modelli, bisogna confrontare il TDC risultante e la stessa coda.

\subsection{Copule archimedee}
Una copula archimedea bivariata è costruita da un generatore $\psi$:
\begin{equation}
  C(u,v)=\psi^{[-1]}\!\left(\psi(u)+\psi(v)\right),
  \label{eq:arch}
\end{equation}
dove $\psi$ è continua, decrescente e convessa, con $\psi(1)=0$.
Il paper esprime i TDC tramite le derivate unilaterali del generatore: la
dipendenza superiore richiede $\psi'(1)=0$; se $\psi'(1)<0$ si ha
indipendenza superiore. La dipendenza inferiore richiede $\psi'(0)=-\infty$
e un generatore stretto; se $\psi'(0)>-\infty$ (o il generatore non è
stretto) si ha indipendenza inferiore.

\subsection{Mappa qualitativa delle famiglie}
\begin{center}
\small
\begin{tabular}{p{2.7cm}p{3.3cm}p{4.4cm}}
\toprule
Famiglia & TDC asintotico & Cosa comunica il modello\\
\midrule
Indipendenza & $\lambda_L=\lambda_U=0$ & Nessun legame estremo; utile come riferimento.\\
Clayton & $\lambda_L>0$, $\lambda_U=0$ (per $\theta>0$) & Legame concentrato nella coda inferiore.\\
Gumbel--Hougaard & $\lambda_U>0$, $\lambda_L=0$ (per $\theta>1$) & Legame concentrato nella coda superiore.\\
Ali--Mikhail--Haq & $\lambda_L=\lambda_U=0$ & Può avere dipendenza centrale senza dipendenza asintotica.\\
Gaussiana & $\lambda_L=\lambda_U=0$ per $|\rho|<1$ & La correlazione può essere alta, ma la coda estrema resta asintoticamente indipendente.\\
$t$-copula & $\lambda_L=\lambda_U>0$ in genere & Code pesanti e dipendenza di coda simmetrica; dipende da correlazione e gradi di libertà.\\
\bottomrule
\end{tabular}
\end{center}

\paragraph{Attenzione al parametro.}
In una $t$-copula, più gradi di libertà significano code meno pesanti e, a
parità di correlazione, un TDC più piccolo; non è corretto confrontare
direttamente questo parametro con $\theta$ di Clayton o Gumbel senza prima
guardare cosa produce sui TDC.

\subsection{Distribuzioni ellittiche}
Una distribuzione ellittica ha, quando esiste, densità
\begin{equation}
 f(x)=|\Sigma|^{-1/2}\,g\!\left((x-\mu)^\top\Sigma^{-1}(x-\mu)\right),
 \qquad x\in\R^n,
 \label{eq:elliptical}
\end{equation}
dove $\mu$ è la posizione, $\Sigma$ la matrice di dispersione e $g$ il
generatore di densità. Le curve di livello sono ellissi, da cui il nome.

Se $g$ è regolarmente variabile con indice $-\alpha/2-1$, allora in
dimensione due la dipendenza superiore e quella inferiore coincidono e
valgono
\begin{equation}
 \lambda_U=\lambda_L=\lambda
 =\frac{\displaystyle\int_0^{h(\rho)}
       \frac{u^\alpha}{\sqrt{1-u^2}}\,du}
       {\displaystyle\int_0^1\frac{u^\alpha}{\sqrt{1-u^2}}\,du},
 \qquad
 h(\rho)=\left(1+\frac{(1-\rho)^2}{1-\rho^2}\right)^{-1/2},
 \label{eq:elliptical-tdc}
\end{equation}
con $\rho=\sigma_{12}/\sqrt{\sigma_{11}\sigma_{22}}$. Un indice $\alpha$
piccolo corrisponde a una coda più pesante e, a parità di $\rho$, a maggiore
probabilità di estremi congiunti. La normale ha coda troppo sottile
($\alpha=\infty$) e nessuna tail dependence; la Student-$t$ ha
$\alpha=\theta$ (gradi di libertà) e TDC positivo.

\section{Come si stima il TDC (Sec.~3.4, pp.~75--78)}

Supponiamo di avere $m$ osservazioni iid bivariate e marginali continue.
Siano $R_1^{(j)},R_2^{(j)}$ i ranghi dell'osservazione $j$ e $k=k(m)$ un numero
di osservazioni estreme, con
\begin{equation}
  k\longrightarrow\infty,\qquad k/m\longrightarrow0.
\end{equation}
\textbf{Nota:} qui $k$ è la lettera latina ``k'' (si legge spesso ``cappa'')
usata come soglia numerica; non è una nuova misura chiamata kappa.

\subsection{Mini-esempio con $m=20$ e $k=4$}
Con $m=20$, i quattro ranghi più alti sono $17,18,19,20$. Immaginiamo che
due osservazioni abbiano contemporaneamente $(R_1,R_2)$ sopra $16$:
\begin{center}
\begin{tabular}{c|c|c|c}
Osservazione & $R_1$ & $R_2$ & entra nella coda upper?\\
\midrule
1 & 20 & 18 & sì\\
2 & 19 & 12 & no\\
3 & 17 & 17 & sì\\
4 & 11 & 20 & no\\
\multicolumn{4}{c}{\small le altre 16 osservazioni hanno almeno un rango $\leq16$}
\end{tabular}
\end{center}
Sono entrati 2 punti nell'intersezione delle due code. Il primo stimatore è
quindi un singolo numero:
\[
  \widehat\lambda_{U}^{(1)}=\frac{2}{k}=\frac{2}{4}=0.5.
\]
Non stiamo stimando una probabilità con tutti i 20 punti: il denominatore $k$
normalizza rispetto alle circa $k$ occasioni in cui una variabile è nella
propria coda.

\begin{figure}[ht]
\centering
\begin{tikzpicture}[x=0.7cm,y=0.7cm,font=\small]
  \draw[->] (0,0)--(5.2,0) node[right] {$R_1$};
  \draw[->] (0,0)--(0,5.2) node[above] {$R_2$};
  \draw[dashed] (3.6,0)--(3.6,5) node[above] {$m-k=16$};
  \draw[dashed] (0,3.6)--(5,3.6) node[right] {$m-k=16$};
  \fill[orange!20] (3.6,3.6) rectangle (5,5);
  \node[align=center] at (4.3,4.45) {intersezione\\$2$ punti};
  \foreach \x/\y in {4.6/4.1,4.1/2.7,3.8/3.8,2.4/4.5,1.6/1.2,2.7/2.4}{\fill[blue!65] (\x,\y) circle (2pt);}
  \fill[red!70] (4.6,4.1) circle (2.7pt);
  \fill[red!70] (3.8,3.8) circle (2.7pt);
\end{tikzpicture}
\caption{Schema dei ranghi per $m=20,k=4$: la zona ombreggiata contiene i punti con entrambi i ranghi sopra 16. Diagramma didattico.}
\label{fig:ranks}
\end{figure}

\subsection{I due stimatori non parametrici}
Il primo conta quante osservazioni cadono contemporaneamente nel quadrato di
coda:
\begin{align}
 \widehat\lambda_{U,m}^{(1)}
   &=\frac1k\sum_{j=1}^m
     \ind\!\left(R_1^{(j)}>m-k,\ R_2^{(j)}>m-k\right),\\
 \widehat\lambda_{L,m}^{(1)}
   &=\frac1k\sum_{j=1}^m
     \ind\!\left(R_1^{(j)}\leq k,\ R_2^{(j)}\leq k\right).
 \label{eq:np1}
\end{align}
Il secondo usa l'unione nella coda superiore:
\begin{equation}
 \widehat\lambda_{U,m}^{(2)}
 =2-\frac1k\sum_{j=1}^m
   \ind\!\left(R_1^{(j)}>m-k\ \text{o}\ R_2^{(j)}>m-k\right).
 \label{eq:np2}
\end{equation}
Nel mini-esempio, ciascuna coda contiene 4 punti e l'intersezione 2; quindi
l'unione contiene $4+4-2=6$ punti e
$\widehat\lambda_U^{(2)}=2-6/4=0.5$. La formula non è magia: sottrae dalla
dimensione teorica dell'unione normalizzata ($2$) la frazione di osservazioni
che appartiene ad almeno una coda.

Il valore di $k$ crea il classico compromesso varianza--distorsione: $k$
piccolo usa pochi dati (stima rumorosa), $k$ grande si allontana dalla coda
(distorsione). Uno stimatore restituisce un singolo numero tra $0$ e $1$,
interpretabile come quota limite stimata, non una nuova serie temporale.

\subsection{Parametrico e non parametrico}
Un metodo non parametrico usa direttamente ranghi e conteggi, senza imporre una
famiglia per la distribuzione congiunta. È più flessibile, ma sensibile alla
scelta di $k$. Nel metodo parametrico per modelli ellittici si stimano invece
l'indice di coda $\alpha$ e il parametro $\rho$, poi si sostituiscono nella
formula \eqref{eq:elliptical-tdc}:
\begin{equation}
 \widehat\lambda_{U,m}^{(3)}
 =\widehat\lambda_{L,m}^{(3)}
 =\lambda\!\left(\widehat\alpha_m,\widehat\rho_m\right).
 \label{eq:param}
\end{equation}
Il metodo di Hill per $\alpha$ è, intuitivamente, un modo per usare solo le
osservazioni più estreme e misurare quanto rapidamente decresce la coda.
Schmidt cita anche stime robuste per $\rho$. L'efficienza aumenta se il
modello ellittico è plausibile, ma una famiglia sbagliata può compromettere la
stima; inoltre la trasformazione integrale può rendere distorto lo stimatore
del TDC anche quando $\widehat\alpha$ è non distorto.

\section{Confronto Monte Carlo (Sec.~3.5, pp.~78--81)}

Vengono generate 1\,000 repliche di campioni iid di dimensione
$m=500,1000,2000$ da una $t$ bivariata con gradi di libertà
$\theta=1.5,2,3$. Qui $\theta=\alpha$ è l'indice di variazione regolare.
Cambiare $m$ studia l'effetto della dimensione campionaria; cambiare $\theta$
non è invece solo un altro esperimento sulla dimensione: confronta scenari di
rischio con code diverse.

\begin{center}
\small
\begin{tabular}{c|c|c|c}
\toprule
$\theta=\alpha$ & Coda & TDC vero (paper) & Lettura qualitativa\\
\midrule
1.5 & molto pesante & $0.2296$ & più estremi congiunti\\
2 & intermedia & $0.1817$ & caso di riferimento\\
3 & meno pesante & $0.1161$ & più vicino alla normale\\
\bottomrule
\end{tabular}
\end{center}

A parità di $\rho$, $\alpha$ piccolo significa coda più pesante e TDC più
grande. Quando $\theta$ cresce la $t$ assomiglia di più alla normale e il TDC
diminuisce. Per ogni scenario, aumentando $m$ diminuiscono bias e RMSE. Nella
simulazione ellittica lo stimatore parametrico $\widehat\lambda^{(3)}$ è il
migliore; per $m=2000,\theta=2$ bias/RMSE (moltiplicati per $10^3$) sono circa
$5.4/17.0$ contro $19.1/41.1$ e $32.9/47.7$ dei due non parametrici. Se la
distribuzione è non ellittica o ignota, il paper raccomanda il primo stimatore
non parametrico.

\section{Rendimenti finanziari: leggere un plateau (Sec.~3.6, pp.~81--84)}

Schmidt analizza rendimenti logaritmici giornalieri trasformati in perdite:
\begin{itemize}[nosep]
  \item $D_1$: BMW e Deutsche Bank, 2\,347 osservazioni (1992--2001);
  \item $D_2$: DEM/USD e JPY/USD, 3\,126 osservazioni (1989--2001).
\end{itemize}

\paragraph{Stima empirica, passo per passo.}
Si trasformano le osservazioni in ranghi o percentili, si fissa una serie di
valori di $k$, si conta quante coppie cadono contemporaneamente nella coda e
si calcola $\widehat\lambda_U^{(1)}(k)$. Poi si disegna la stima contro $k$.
Un plateau a piccoli $k$ suggerisce che, in una regione di soglie, la quota
stimata è relativamente stabile: è l'indizio grafico di dipendenza di coda
usato nel paper. Non è una prova automatica e non elimina la scelta soggettiva
di $k$.

Per $D_1$ il paper sceglie $k$ tra 80 e 110 e ottiene circa
$\widehat\lambda_U=0.31$; per $D_2$ sceglie $k$ tra 40 e 90 e ottiene circa
$\widehat\lambda_U=0.17$. Con $k$ troppo piccolo la curva oscilla per la
varianza; con $k$ troppo grande incorpora osservazioni non più veramente
estreme e cresce la distorsione.

\section{VaR: perché la coda cambia la misura di rischio (Sec.~3.7, pp.~84--88)}

Il VaR a livello $p$ è un quantile della perdita: al livello $1\%$ è una
soglia tale che circa l'1\% delle perdite è ancora più grande. Per un
portafoglio, una grande perdita congiunta di due titoli può spostare molto
questo quantile. Perciò contano sia le code marginali sia la copula.

\begin{figure}[ht]
\centering
\begin{tikzpicture}[x=0.9cm,y=0.75cm,font=\small]
  \draw[->] (0,0)--(8.4,0) node[right] {perdita};
  \draw[->] (0,0)--(0,4.2) node[above] {densità};
  \draw[domain=0.2:7.8,smooth,thick,blue!70] plot (\x,{3.4*exp(-0.5*((\x-3.0)/1.25)^2)});
  \fill[red!25] (0,0) -- plot[domain=0.2:0.75,smooth] (\x,{3.4*exp(-0.5*((\x-3.0)/1.25)^2)}) -- (0.75,0)--cycle;
  \draw[dashed,red!70,thick] (0.75,0)--(0.75,1.0) node[above right] {$q_{1\%}$};
  \node[red!70,align=center] at (0.45,2.0) {coda\\1\%};
  \node[align=center] at (4.8,3.7) {un quantile molto estremo\\è sensibile alla forma della coda};
\end{tikzpicture}
\caption{Schema concettuale di una distribuzione di perdite. La coda evidenziata rappresenta gli esiti oltre il quantile 1\%; non è una stima dei dati del paper.}
\label{fig:var-tail}
\end{figure}

Il terzo campione, $D_3$, contiene FFR/USD e DEM/USD (5\,189 osservazioni,
1984--2002) e non appare ellittico. Per $D_1,D_2,D_3$ Schmidt usa marginali
$t$ e confronta: (i) normale bivariata, (ii) $t$ bivariata, (iii) copula-$t$
con marginali-$t$. I parametri di copula e marginali sono stimati
separatamente con IFM (Inference Functions for Margins).

\paragraph{Che cosa cambia.}
La normale ha code sottili e TDC nullo: può sottostimare VaR a 1\% e 2.5\%.
La $t$ bivariata ha code pesanti e TDC positivo, ma impone che l'intera
distribuzione sia ellittica. La copula-$t$ con marginali separate conserva
code marginali diverse e consente alla dipendenza estrema di avere una forma
propria. È particolarmente utile per $D_3$: il grado di libertà della copula è
circa $1.11$, mentre quelli delle marginali sono circa $4.63$ e $5.15$.

I risultati del paper mostrano che $D_1$ e $D_2$ sono abbastanza vicini a una
struttura ellittica, mentre per $D_3$ separare copula e marginali migliora la
stima. La copula può quindi cambiare materialmente il VaR anche quando le
marginali sono tenute fisse: modifica la probabilità di osservare perdite
congiunte nella regione estrema.

\section{Riepilogo operativo per la discussione}

\subsection*{Cosa dovrei aver capito}
\begin{enumerate}[leftmargin=*]
  \item La correlazione descrive una relazione media; il TDC chiede cosa
  succede quando la soglia tende all'estremo.
  \item $\lambda_U$ e $\lambda_L$ sono probabilità condizionate limite,
  rispettivamente in alto a destra e in basso a sinistra.
  \item Una copula separa marginali e dipendenza: il TDC è una proprietà della
  copula per marginali continue.
  \item Normale e $t$ possono avere dipendenza centrale simile ma comportamento
  di coda opposto; una $t$-copula può essere utile quando la coda è importante.
  \item Lo stimatore non parametrico conta ranghi estremi e richiede $k$; il
  parametrico usa una famiglia e stima $\alpha,\rho$.
  \item Nel VaR la scelta di code e copula può spostare la stima del rischio.
\end{enumerate}

\subsection*{Cosa è ancora ostico e cosa rileggere}
La condizione di variazione regolare e la formula integrale
\eqref{eq:elliptical-tdc} sono il passaggio matematico più tecnico: rileggere
pp.~72--74. Il secondo stimatore non parametrico e il motivo per cui usa
l'unione sono da rileggere a p.~77. Per la parte empirica, rileggere il
compromesso varianza--distorsione nelle pp.~82--83 e il confronto dei tre
modelli VaR nelle pp.~85--88.

\subsection*{Parole chiave}
tail dependence, TDC, $\lambda_U$, $\lambda_L$, indipendenza asintotica,
copula, marginali, ranghi, variazione regolare, indice $\alpha$, copula-
Gaussiana, $t$-copula, Hill estimator, bias--varianza, VaR, IFM.

\subsection*{Domande da portare al relatore}
\begin{description}[leftmargin=3.1cm,style=nextline]
  \item[Concettuali] Qual è il rapporto tra TDC e intermediate dependence?
  Per la tesi servono entrambe le code o soprattutto la lower tail? Quanto è
  importante distinguere indipendenza asintotica da indipendenza tout court?
  \item[Metodologiche] Come scegliere $k$ in modo replicabile? Quando è
  ragionevole imporre una famiglia ellittica? È preferibile stimare una copula
  dopo aver calibrato le marginali o usare una distribuzione congiunta?
  \item[Empiriche] Come controllare l'ipotesi iid sui rendimenti? Quale
  diagnostica usare oltre al plateau? Come misurare l'incertezza della stima
  del TDC e del VaR?
  \item[Collegamenti] Hua--Joe sembra spostare l'attenzione dalla distinzione
  binaria dipendenza/\allowbreak indipendenza verso la velocità di decadimento;
  Durante
  \emph{et al.} sembra collegare la diagonal section e gli ordini di coda a
  una descrizione più generale. Conviene considerarli estensioni teoriche del
  quadro introduttivo di Schmidt o approcci alternativi?
\end{description}

\section{Legenda della notazione e della pronuncia}
\label{sec:legenda}

Questa tabella serve per rileggere le formule senza dover ricordare tutto a
memoria.

\begin{center}
\small
\begin{tabular}{p{0.18\textwidth}p{0.27\textwidth}p{0.45\textwidth}}
\toprule
Simbolo & Come si legge & Significato semplice\\
\midrule
$X,X_1,X_2$ & ``ics'' & variabili aleatorie, ad esempio rendimenti o perdite\\
$F,F_1,F_2$ & ``effe'' & funzioni di distribuzione congiunta o marginali\\
$F_i^{-1}(v)$ & ``effe i inversa'' & quantile, cioè soglia al percentile $v$\\
$\lambda$ & ``lambda'' & coefficiente di dipendenza di coda\\
$\lambda_U,\lambda_L$ & ``lambda upper/lower'' & TDC superiore/inferiore\\
$\rho$ & ``rho'' & correlazione o parametro di dipendenza ellittico\\
$\theta$ & ``teta'' & parametro di famiglia o gradi di libertà della $t$\\
$\alpha$ & ``alfa'' & indice della pesantezza della coda\\
$\psi$ & ``psi'' & generatore di una copula archimedea\\
$\psi^{[-1]}$ & ``psi pseudo-inversa'' & pseudo-inversa del generatore, non inversa della normale\\
$\Sigma$ & ``sigma maiuscola'' & matrice di scala/dispersione\\
$\mu$ & ``mu'' & vettore di posizione o media\\
$u,v$ & ``u, vi'' & percentili in $[0,1]$, non rendimenti\\
$k$ & ``cappa'' (lettera latina) & numero di osservazioni estreme usate nella stima\\
$m$ & ``emme'' & dimensione del campione\\
$\Prob$ o $P$ & ``probabilità'' & probabilità di un evento\\
$\lim$ & ``limite'' & valore verso cui tende una quantità\\
$\ind(\cdot)$ & ``indicatrice'' & vale 1 se la condizione è vera, 0 altrimenti\\
$t$ & ``t'' & distribuzione di Student, con code più pesanti della normale\\
$\operatorname{VaR}$ & ``vi-a-erre'' & quantile estremo della perdita\\
$\operatorname{RMSE}$ & ``erre-emme-esse-e'' & errore che combina bias e variabilità\\
\bottomrule
\end{tabular}
\end{center}

\paragraph{Altri simboli ricorrenti.}
$v\uparrow1$ significa che $v$ si avvicina a 1 da sinistra; $v\downarrow0$
che si avvicina a 0 da destra. Il simbolo $\mid$ nella probabilità
condizionata si legge ``dato che''. Il simbolo $\top$ indica la trasposta di
un vettore; $|\Sigma|$ è il determinante della matrice e $\Sigma^{-1}$ la sua
inversa.

\end{document}
